% Part: incompleteness
% Chapter: representability-in-q
% Section: prim-rec

\documentclass[../../../include/open-logic-section]{subfiles}

\begin{document}

\olfileid{inc}{req}{pri}
\olsection{ఆదిమ పునరావృత్తిని అనుకరించడం}

ఇప్పుడు బీటా ప్రమేయాన్ని వాడి, ఆదిమ పునరావృత్త నిర్వచనాన్ని సక్రమ
కనిష్ఠీకరణతో ``అనుకరించవచ్చని'' చూపవచ్చు. $f(\vec x)$,
$g(\vec x,y,z)$లు ఉన్నాయని అనుకుందాం. అప్పుడు $f$, $g$ల నుంచి ఆదిమ
పునరావృత్తి ద్వారా నిర్వచించిన ప్రమేయం~$h(\vec x,y)$:
\begin{align*}
h(\vec x, 0) & =  f(\vec x) \\
h(\vec x, y+1) & =  g(\vec x, y, h(\vec x, y)).
\end{align*}
\sourcecorrection{OLTEREQ-003}{ఆంగ్ల మూలం పరిచయ వాక్యంలో $h(x,\vec z)$ అని చెప్పి, వెంటనే నిర్వచన సమీకరణల్లో $h(\vec x,y)$ను వాడింది. ఆ సమీకరణాలు, తరువాతి నిరూపణ నిర్దేశించే $h(\vec x,y)$నే వాడాం.}
ప్రాథమిక ప్రమేయాలనూ, వాటి నుంచి సంయుక్తం, సక్రమ కనిష్ఠీకరణలతో నిర్వచించిన
$\beta$ వంటి ప్రమేయాలనూ వాడుతూ, సంయుక్తం, సక్రమ కనిష్ఠీకరణ మాత్రమే ఉపయోగించి
$f$, $g$ల నుంచి~$h$ను నిర్వచించవచ్చని చూపాలి.

\begin{lem}
\ollabel{lem:prim-rec}
$f$, $g$ల నుంచి ఆదిమ పునరావృత్తితో~$h$ను నిర్వచించగలిగితే, $f$, $g$,
$\Zero$, $\Succ$, $\Proj{n}{i}$, $\Add$, $\Mult$, $\Char{=}$ ప్రమేయాల నుంచి
సంయుక్తం, సక్రమ కనిష్ఠీకరణలతో దాన్ని నిర్వచించవచ్చు.
\end{lem}

\begin{proof}
మొదట $\hat h(\vec x,y)$ అనే సహాయక ప్రమేయాన్ని నిర్వచించండి. అది కనిష్ఠ
సంఖ్య~$d$ను ఇస్తుంది; ఆ~$d$ కింది షరతులు నెరవేర్చే క్రమాన్ని సంకేతీకరిస్తుంది:
\begin{enumerate}
\item $(d)_0=f(\vec x)$, మరియు
\item ప్రతి $i<y$కూ $(d)_{i+1}=g(\vec x,i,(d)_i)$,
\end{enumerate}
ఇక్కడ ఇప్పుడు $(d)_i$ అనేది $\beta(d,i)$కు సంక్షిప్త రూపం. మరో మాటలో,
$\hat h$ అనేది $\tuple{h(\vec x,0),h(\vec x,1),\dots,h(\vec x,y)}$ క్రమాన్ని
ఇస్తుంది. $\hat h$ను ఇలా వ్రాయవచ్చు:
\[
\hat h(\vec x, y) = \umin{d}{(\beta(d,0) = f(\vec x) \land \bforall{i <
  y}{\beta(d,i+1) = g(\vec x, i,\beta(d,i)})}.
\]
ఇక్కడ ఆదిమ పునరావృత్తి అవసరం లేదని గమనించండి; కనిష్ఠీకరణ మాత్రమే చాలు.
బీటా ప్రమేయ ఉపపత్తి~\olref[bet]{lem:beta} వల్ల మనం కనిష్ఠీకరించే ప్రమేయం
సక్రమమైనది.

కానీ ఇప్పుడు
\[
h(\vec x, y) = \beta(\hat h(\vec x, y), y),
\]
కాబట్టి, ప్రాథమిక ప్రమేయాల నుంచి సంయుక్తం, సక్రమ కనిష్ఠీకరణ మాత్రమే వాడి
$h$ను నిర్వచించవచ్చు.
\end{proof}

\end{document}
